\documentclass[12pt,a4paper,onecolumn]{article}
\usepackage[margin=27mm]{geometry}
\usepackage{fontspec,xeCJK,amsmath,amssymb,booktabs}
\usepackage[hidelinks]{hyperref}
\setmainfont[Path=C:/Windows/Fonts/,BoldFont=timesbd.ttf,ItalicFont=timesi.ttf]{times.ttf}
\setmonofont[Path=C:/Windows/Fonts/]{times.ttf}
\setCJKmainfont[Path=C:/Windows/Fonts/,BoldFont=kaiu.ttf]{kaiu.ttf}
\renewcommand{\normalsize}{\fontsize{12bp}{19bp}\selectfont}
\setlength{\parindent}{2em}
\setlength{\parskip}{4bp}
\setlength{\emergencystretch}{2em}
\widowpenalty=10000
\clubpenalty=10000
\renewcommand{\contentsname}{目錄}
\renewcommand{\refname}{參考資料}
\hypersetup{pdftitle={盲式數位浮水印的 JPEG 抗壓縮分析與 SVD 區塊改良},pdfauthor={林易辰}}
\begin{document}
\normalsize
\begin{titlepage}
\centering
\vspace*{15mm}
{\Large 矩陣理論及應用\quad 期末專題\par}
\vspace{10mm}
{\LARGE 盲式數位浮水印的 JPEG 抗壓縮分析\\[3mm]與 SVD 區塊改良\par}
\vspace{8mm}
{\large 內容審閱稿\quad 最終方法與結果待定\par}
\vspace{8mm}
\begin{tabular}{ll}
作者 & 林易辰\\
學號 & B143012007\\
學校／系級 & ［待填］\\
授課教師 & ［待填］\\
繳交日期 & ［待填］
\end{tabular}

\vspace{10mm}
\begin{minipage}{0.91\textwidth}
\section*{摘要}
同一張含浮水印影像經 JPEG 壓縮後，外觀看起來仍可能相似，擷取的浮水印位元卻可能出錯。本專題以 \texttt{blind\_watermark} 的 DWT–DCT–SVD 盲擷取流程為對象，說明原版如何將位元寫入區塊的前兩個奇異值，以及如何以餘數門檻讀回。本研究以 SVD 區塊為改良範圍，在預定的 JPEG 攻擊條件下配對比較原版與候選，量測逐位錯誤率及影像品質。最終候選與實測結論待驗證後填入。

\noindent\textbf{關鍵字}\quad 盲式數位浮水印、JPEG、奇異值分解、區塊改良、位元錯誤率
\end{minipage}
\end{titlepage}
\setcounter{page}{2}

\section{緒論}

\subsection{研究動機}
將照片寄出或上傳後，平台可能重新儲存為 JPEG。觀看者仍可辨認原來的照片，但每個像素的值不一定與上傳前相同。若照片裡另有一段用來識別來源的浮水印，判斷浮水印是否保留下來，需要重新讀取並比對位元。單憑視覺相似度，無法知道其中有幾位被讀錯。

本研究使用開放原始碼 \texttt{blind\_watermark} 的固定版本 \cite{core,interface}。原版以矩陣分解取得區塊的奇異值，藉由改變奇異值的餘數寫入位元。擷取時只看餘數落在哪一側，便把該奇異值的判斷化成 0 或 1。假如兩個餘數都落在同一側，其中一個緊鄰分界，另一個靠近嵌入時的理想位置，原版仍給出相同的判斷值。這正是本專題選擇修改的步驟。

\subsection{研究問題與範圍}
只修改 \texttt{blind\_watermark} 的 SVD 區塊，能否在可比的影像品質下，降低 JPEG 壓縮後的盲擷取位元錯誤率？研究以固定版本的原版實作為基準，JPEG 作黑盒攻擊。各組比較固定輸入影像、訊息、密鑰及攻擊設定，並分別量測擷取錯誤和影像失真。

W0 建立原版重現流程，W0-D 提供固定攻擊影像與判決診斷 \cite{w0,w0d}。最終候選由前置診斷與驗證決定。正文以選定方法和可核對的配對證據為主。

\section{數位浮水印、JPEG 與 SVD 基礎}

\subsection{嵌入與盲擷取}
數位浮水印將短訊息寫入影像。本文把訊息表示為一串 0 和 1。設原圖為 \(I\)，訊息為 \(b=(b_1,\ldots,b_L)\)，嵌入後的影像為 \(I_w\)。JPEG 編解碼後得到 \(\widehat I_{w,Q}\)，其中 \(Q\) 是編碼器的品質設定。擷取器由受測影像、訊息長度、固定參數及密鑰取得估計位元 \(\widehat b\)。原圖及正確位元留在評估端，所以擷取屬於盲擷取。

\[
 (I,b)\ \longrightarrow\ I_w
 \ \xrightarrow{\text{JPEG 編解碼}}\ \widehat I_{w,Q}
 \ \xrightarrow{\text{盲擷取}}\ \widehat b .
\]
無攻擊對照直接由存檔的 \(I_w\) 擷取。它用來檢查寫入、儲存、重讀及擷取的整條流程。若這一步已失敗，需先檢查基準，才能解讀 JPEG 攻擊下的差異。

\subsection{JPEG 為何可能改變判決}
本研究的 JPEG 是有損影像編碼。編碼器將影像轉為可壓縮表示，量化會捨去部分精度，解碼後的像素因此可能改變 \cite{jpeg}。JPEG 品質設定可用來建立不同攻擊條件，實驗仍須記錄實際的編碼器、色度取樣及其他設定。本文把 JPEG 視為黑盒，只控制輸入和編解碼條件，並讀回實際解碼影像。

浮水印程式本身也使用 DCT，但用途不同。它先在小區塊內整理資料，再從變換後的矩陣取得奇異值以讀寫浮水印。JPEG 改變影像像素後，擷取器重算的區塊矩陣及其奇異值可能改變。若餘數跨過判決邊界，部分區塊的票可能變成另一個位元。

\subsection{位元與影像品質怎麼量}
位元錯誤率（bit error rate，BER）是逐一比較位置所得的比例。以 \(\mathbf 1[\cdot]\) 表示條件成立時取 1，否則取 0，則
\[
 \mathrm{BER}=\frac{1}{L}\sum_{i=1}^{L}\mathbf 1[\widehat b_i\ne b_i].
\]
例如讀回 32 位而其中 2 位錯誤，BER 為 \(2/32=6.25\%\)。計算前先確認長度和位元順序，報告時同時保留錯誤數與總位元數。

影像品質以峰值訊雜比（peak signal-to-noise ratio，PSNR）描述。對同尺寸的 8 位元 RGB 影像 \(X,Y\)，令 \(N\) 為三色通道的樣本總數，則
\[
 \mathrm{MSE}(X,Y)=\frac{1}{N}\sum_{j=1}^{N}(X_j-Y_j)^2,\qquad
 \mathrm{PSNR}(X,Y)=10\log_{10}\frac{255^2}{\mathrm{MSE}(X,Y)}.
\]
均方誤差為零時，PSNR 記為 \(+\infty\)。\(\mathrm{PSNR}(I,I_w)\) 描述嵌入失真，\(\mathrm{PSNR}(I,\widehat I_{w,Q})\) 描述嵌入與 JPEG 的合併失真。兩者用不同欄位記錄。若兩種擷取器共用同一份受測影像，同條件的兩組 PSNR 相同。

\subsection{奇異值分解基礎}

\subsubsection{一個區塊如何寫成矩陣成分}
用 \(B\) 表示一個送進 SVD 的實數 \(4\times4\) 矩陣。其奇異值分解可寫成
\[
 B=U\Sigma V^{\mathsf T}
  =\sum_{j=1}^{4}\sigma_j u_jv_j^{\mathsf T},\qquad
 \sigma_1\ge \sigma_2\ge \sigma_3\ge \sigma_4\ge0.
\]
\(U\) 與 \(V\) 的直行向量（columns）各自正交歸一，\(\Sigma\) 的對角元素是奇異值。以 \(B^{\mathsf T}B\) 來看，\(\sigma_j\) 是其特徵值的非負平方根。外積 \(u_jv_j^{\mathsf T}\) 給出一個秩一成分的形狀，\(\sigma_j\) 控制該成分的大小。

以 \(B=\operatorname{diag}(5,2)\) 為小型例子。兩個方向分別對應左上角與右下角。把第一奇異值由 5 改為 6，便在第一個方向上增加 1，而第二方向保持原值。實際區塊的方向通常並非座標軸，這個改量會分布在矩陣的多個位置。

\subsubsection{修改奇異值如何改變區塊}
保留 \(U,V\)，將前兩項改為 \(\sigma'_1,\sigma'_2\) 時，重建的矩陣 \(B'\) 與 \(B\) 相差
\[
 B'-B=(\sigma'_1-\sigma_1)u_1v_1^{\mathsf T}
     +(\sigma'_2-\sigma_2)u_2v_2^{\mathsf T}.
\]
兩個外積的 Frobenius 內積為
\((u_1^{\mathsf T}u_2)(v_1^{\mathsf T}v_2)=0\)，各自平方範數為 1。因此
\[
 \|B'-B\|_F^2=(\sigma'_1-\sigma_1)^2+(\sigma'_2-\sigma_2)^2.
\]
這是固定方向下的區塊矩陣等式。經過逆變換、像素量化、裁切及 JPEG 後，整張圖的失真和位元錯誤率仍應從實際影像計算。

\section{原版 DWT–DCT–SVD 演算法與實驗方法}

\subsection{從彩色影像到嵌入區塊}
原版 common 模式將影像轉成 YUV 三通道，對每個通道做 Haar 離散小波轉換（DWT），取低頻近似部分，再切成 \(4\times4\) 區塊。各塊做離散餘弦轉換（DCT），把 16 個係數展開並按密鑰產生的順序重排，還原成 \(4\times4\) 矩陣後做 SVD。位元會在通道和區塊間重複寫入，以便在擷取時合併多次觀察 \cite{core}。

程式把前兩個奇異值寫成 \texttt{s[0]}、\texttt{s[1]}。本文用 \(j=1,2\) 對應，原版步長為 \(d_1=36,\ d_2=20\)。令要寫入該塊的位元為 \(b\in\{0,1\}\)，則
\[
 \sigma'_j=d_j\left(\left\lfloor\frac{\sigma_j}{d_j}\right\rfloor
                  +\frac14+\frac{b}{2}\right),\qquad j=1,2.
\]
整數部分保留原本所在的步長區間，小數位置則由位元決定。寫入 0 時理想餘數比例為 \(1/4\)，寫入 1 時為 \(3/4\)。例如 \(\sigma_1=50\)、\(d_1=36\)，寫入 0 得到 45，寫入 1 得到 63，餘數分別為 9 和 27。改值後重建區塊，依序還原係數順序、執行逆 DCT 和逆 DWT，最後還原彩色影像。

\subsection{原版擷取如何把區塊判斷合成一位}
擷取端重做相同區塊處理及係數重排。以 \(s_j\) 表示受測影像重算的第 \(j\) 個奇異值，令
\[
 r_j=\frac{s_j\bmod d_j}{d_j}\in[0,1),\qquad
 h(r_j)=\mathbf 1[r_j>1/2].
\]
因此恰等於 \(1/2\) 時判為 0。原版把兩個硬判決以 \(3:1\) 合成區塊分數 \((3h(r_1)+h(r_2))/4\)。這個分數只可能取 \(0,\ 1/4,\ 3/4,\ 1\)。

對應同一訊息位置的區塊分數，跨重複位置與三通道取平均。公開的 \texttt{mode="bit"} 接著對整串平均分數做一維兩群 k-means，把高分群判為 1，再依密鑰還原位元順序 \cite{core,interface}。所以 \(1/2\) 是單個奇異值的餘數門檻。最後的位元判決由所有平均分數共同決定。

\subsection{實驗材料與配對方式}
原版基準使用 W0 的三張固定 RGB 影像、32 位元訊息與密鑰，並保存含水印 PNG 和輸入雜湊 \cite{w0}。W0-D 從每張含水印 PNG 獨立編碼 JPEG，提供 Q10、Q20、Q30 的 4:4:4 封存條件 \cite{w0d}。這些檔案可供比較與診斷。最終候選的完整測試條件須在正式評估前固定，並記錄各條件用途、編碼器版本、品質及色度取樣。

每張影像先作無攻擊 PNG 對照。若候選只修改擷取，原版與候選讀取同一份 PNG 和 JPEG 解碼像素。若候選修改區塊嵌入，兩組從同一原圖出發，使用同一訊息、密鑰及 JPEG 設定，各自從含水印 PNG 獨立編碼。擷取過程只使用受測影像及預先固定的設定。原圖用於計算 PSNR，正確位元用於事後評估。

\subsection{結果如何比較}
逐影像、逐條件保存擷取位元、錯誤數、總位元數及兩種 PSNR。以
\[
 \Delta\mathrm{BER}=\mathrm{BER}_{\mathrm{candidate}}-\mathrm{BER}_{\mathrm{original}}
\]
表示配對差值。負值代表該條件的錯誤率下降。若兩組的含水印影像不同，先依事前固定的品質判準比較嵌入 PSNR，再判讀 BER。三張圖使用同一段 32 位元訊息時，每個條件合計 96 位元。結果只適用於實測樣本與編碼設定。

\section{SVD 區塊改良方法}
本章保留待定。待 G5 前置診斷及最終候選確定後，說明改動的區塊步驟、選擇依據、盲擷取所需條件，以及與原版的差異。

\section{實驗結果與結論}
本章保留待填。最終報告依預定測試條件完整呈現原版與候選的 BER、PSNR、差值與限制。直接影響方法選擇、比較公平性或結論的負面證據一併交代。若實測未支持改善，結論依實際觀察撰寫。

\begin{thebibliography}{9}
\bibitem{core}
Guofei，\texttt{blind\_watermark}，固定提交的
\href{https://github.com/guofei9987/blind_watermark/blob/322c90dc02d45644692095b34d3a6b735b6edbd8/blind_watermark/bwm_core.py}{\texttt{bwm\_core.py}}。
\bibitem{interface}
Guofei，\texttt{blind\_watermark}，同一提交的
\href{https://github.com/guofei9987/blind_watermark/blob/322c90dc02d45644692095b34d3a6b735b6edbd8/blind_watermark/blind_watermark.py}{公開擷取介面}。
\bibitem{jpeg}
ITU-T，\href{https://www.itu.int/rec/T-REC-T.81/en}{Recommendation T.81，JPEG 編碼標準}。
\bibitem{w0}
專題倉庫，\href{https://github.com/Rookie-YCLin/matrix-jpeg-research/blob/f5a47fc1ba40cc35f5794648b4440cb1fed782d1/outputs/run-w0/manifest.json}{W0 基準 manifest}。
\bibitem{w0d}
專題倉庫，\href{https://github.com/Rookie-YCLin/matrix-jpeg-research/blob/168e5070cacbdeeaefab58ed469dd7e90db4aa3f/watermark_baseline/STRESS.md}{W0-D JPEG 壓力診斷}。
\end{thebibliography}
\end{document}
